% Mathématiques 5e — cours V3
% Triangles et parallélogrammes

\section{Objectifs}

\begin{itemize}
\item Connaître et démontrer la somme des angles d'un triangle.
\item Construire un triangle à partir de données partielles.
\item Utiliser les médiatrices et le cercle circonscrit dans des triangles particuliers.
\item Calculer l'aire d'un triangle.
\item Définir et tracer les hauteurs d'un triangle.
\item Savoir que les hauteurs sont concourantes.
\item Définir et tracer les médianes.
\item Démontrer qu'une médiane partage un triangle en deux triangles de même aire.
\item Définir, construire et reconnaître un parallélogramme.
\item Utiliser les propriétés caractéristiques des côtés et des diagonales.
\item Définir rectangle, losange et carré comme parallélogrammes particuliers.
\item Calculer l'aire d'un parallélogramme et de figures composées.
\end{itemize}

\section{1. Somme des angles d'un triangle}

Dans tout triangle :
\[
\boxed{\widehat A+\widehat B+\widehat C=180^\circ}.
\]

Cette propriété est indépendante de la forme du triangle.

\begin{exemple}
Dans un triangle, deux angles mesurent :
\[
48^\circ
\]
et :
\[
67^\circ.
\]

Le troisième angle vaut :
\[
180^\circ-48^\circ-67^\circ=65^\circ.
\]

Donc :
\[
\boxed{65^\circ}.
\]
\end{exemple}

\section{2. Idée d'une démonstration}

On peut tracer par un sommet une droite parallèle au côté opposé.

Les égalités d'angles alternes-internes permettent de reporter les deux autres angles du triangle autour du sommet.

On obtient alors un angle plat :
\[
180^\circ.
\]

Ainsi, les trois angles du triangle ont une somme égale à :
\[
180^\circ.
\]

\begin{remarque}
Au cycle 4, il faut distinguer une mesure observée sur un dessin d'une propriété démontrée.
\end{remarque}

\section{3. Construire un triangle}

Un triangle peut être déterminé à partir de différentes données :
\begin{itemize}
\item trois longueurs ;
\item deux longueurs et l'angle compris ;
\item une longueur et deux angles.
\end{itemize}

La construction utilise règle, compas et rapporteur.

\section{4. Construction avec trois côtés}

On veut construire :
\[
AB=6\ \text{cm},\qquad
AC=5\ \text{cm},\qquad
BC=4\ \text{cm}.
\]

On trace $[AB]$.

Puis :
\begin{itemize}
\item un cercle de centre $A$ et rayon $5\ \text{cm}$ ;
\item un cercle de centre $B$ et rayon $4\ \text{cm}$.
\end{itemize}

Leur point d'intersection donne $C$.

\section{5. Construction avec deux côtés et un angle}

On connaît :
\[
AB=7\ \text{cm},
\]
\[
AC=4\ \text{cm},
\]
et :
\[
\widehat{BAC}=55^\circ.
\]

On trace $[AB]$, puis un angle de :
\[
55^\circ
\]
en $A$.

Sur la seconde demi-droite, on reporte :
\[
AC=4\ \text{cm}.
\]

\section{6. Construction avec un côté et deux angles}

Si :
\[
AB=8\ \text{cm},
\]
\[
\widehat{A}=40^\circ,
\]
et :
\[
\widehat{B}=65^\circ,
\]
on trace d'abord $[AB]$.

On construit ensuite les deux demi-droites correspondant aux angles.

Leur intersection détermine le troisième sommet.

Le troisième angle doit vérifier :
\[
180^\circ-40^\circ-65^\circ=75^\circ.
\]

\section{7. Médiatrices et cercle circonscrit}

La médiatrice d'un segment est la droite perpendiculaire au segment passant par son milieu.

Tout point de la médiatrice de $[AB]$ est équidistant de $A$ et de $B$.

Les trois médiatrices d'un triangle sont concourantes.

Leur point commun est le centre du cercle circonscrit au triangle.

\section{8. Triangles particuliers et cercle circonscrit}

Dans un triangle rectangle, le centre du cercle circonscrit est le milieu de l'hypoténuse.

Dans un triangle isocèle, la médiatrice de la base est aussi un axe de symétrie du triangle.

Ces situations permettent de simplifier les constructions.

\section{9. Aire d'un triangle}

Si $b$ est une base et $h$ la hauteur associée :
\[
\boxed{A=\dfrac{b\times h}{2}}.
\]

\begin{exemple}
Pour :
\[
b=9\ \text{cm},
\qquad
h=4\ \text{cm},
\]
on obtient :
\[
A=\dfrac{9\times4}{2}=18\ \text{cm}^2.
\]
\end{exemple}

\section{10. Hauteur d'un triangle}

\begin{definition}
Une hauteur d'un triangle est une droite passant par un sommet et perpendiculaire au côté opposé, ou à la droite qui le porte.
\end{definition}

\coursefigure{/images/mathematiques/5eme/triangle-droites.svg}{Triangle avec une hauteur, une médiane et une médiatrice représentées différemment.}{Hauteur, médiane et médiatrice ne répondent pas à la même définition et ne doivent pas être confondues.}

Dans un triangle obtus, une hauteur peut se situer à l'extérieur du triangle.

\section{11. Orthocentre}

Les trois hauteurs d'un triangle se coupent en un même point.

On dit qu'elles sont concourantes.

Le point de concours s'appelle l'orthocentre.

\begin{propriete}
Les trois hauteurs d'un triangle sont concourantes.
\end{propriete}

\section{12. Médiane}

\begin{definition}
Une médiane d'un triangle est une droite passant par un sommet et par le milieu du côté opposé.
\end{definition}

Chaque triangle possède trois médianes.

\section{13. Médiane et aire}

Considérons un triangle $ABC$ et $M$, milieu de $[BC]$.

Alors :
\[
BM=MC.
\]

Les triangles $ABM$ et $AMC$ ont la même hauteur issue de $A$.

Leur base a la même longueur.

Donc :
\[
A_{ABM}=A_{AMC}.
\]

\begin{propriete}
Une médiane partage un triangle en deux triangles de même aire.
\end{propriete}

\section{14. Triangles particuliers}

Dans un triangle isocèle, certaines droites remarquables peuvent coïncider.

Par exemple, dans un triangle isocèle en $A$, la droite issue de $A$ passant par le milieu de la base est à la fois :
\begin{itemize}
\item médiane ;
\item hauteur ;
\item médiatrice de la base ;
\item axe de symétrie.
\end{itemize}

Dans un triangle équilatéral, cette coïncidence se produit depuis chacun des sommets.

\section{15. Parallélogramme}

\begin{definition}
Un parallélogramme est un quadrilatère dont les côtés opposés sont parallèles deux à deux.
\end{definition}

Si $ABCD$ est un parallélogramme :
\[
(AB)\parallel(CD)
\]
et :
\[
(BC)\parallel(AD).
\]

\section{16. Construire un parallélogramme}

Pour construire $ABCD$ à partir de trois sommets :
\begin{enumerate}
\item tracer par $C$ la parallèle à $(AB)$ ;
\item tracer par $A$ la parallèle à $(BC)$ ;
\item leur intersection donne $D$.
\end{enumerate}

On peut aussi utiliser les diagonales et leur propriété caractéristique.

\section{17. Côtés opposés}

Dans un parallélogramme, les côtés opposés ont la même longueur :
\[
AB=CD,
\]
et :
\[
BC=AD.
\]

Cette propriété permet de calculer des longueurs ou de reconnaître un parallélogramme dans certaines configurations.

\section{18. Diagonales}

\coursefigure{/images/mathematiques/5eme/parallelogramme.svg}{Parallélogramme avec diagonales se coupant en leur milieu et une hauteur relative à une base.}{Les diagonales d'un parallélogramme se coupent en leur milieu ; cette propriété est caractéristique.}

\begin{propriete}
Les diagonales d'un parallélogramme se coupent en leur milieu.
\end{propriete}

Réciproquement, si les diagonales d'un quadrilatère se coupent en leur milieu, alors ce quadrilatère est un parallélogramme.

\section{19. Propriété caractéristique par les côtés}

Si un quadrilatère possède deux côtés opposés parallèles et de même longueur, alors c'est un parallélogramme.

De même, si ses côtés opposés sont de même longueur deux à deux, cela permet aussi de caractériser un parallélogramme.

\section{20. Rectangle}

\begin{definition}
Un rectangle est un parallélogramme qui possède un angle droit.
\end{definition}

Ses quatre angles sont droits.

Ses diagonales :
\begin{itemize}
\item se coupent en leur milieu ;
\item ont la même longueur.
\end{itemize}

\section{21. Losange}

\begin{definition}
Un losange est un parallélogramme dont les quatre côtés ont la même longueur.
\end{definition}

Ses diagonales :
\begin{itemize}
\item se coupent en leur milieu ;
\item sont perpendiculaires.
\end{itemize}

\section{22. Carré}

\begin{definition}
Un carré est à la fois un rectangle et un losange.
\end{definition}

Il possède donc :
\begin{itemize}
\item quatre côtés égaux ;
\item quatre angles droits ;
\item des diagonales de même longueur ;
\item des diagonales perpendiculaires ;
\item des diagonales qui se coupent en leur milieu.
\end{itemize}

\section{23. Aire d'un parallélogramme}

Si $b$ est la longueur d'une base et $h$ la hauteur correspondante :
\[
\boxed{A=b\times h}.
\]

\begin{exemple}
Pour :
\[
b=8\ \text{cm},
\qquad
h=5\ \text{cm},
\]
on obtient :
\[
A=8\times5=40\ \text{cm}^2.
\]
\end{exemple}

\section{24. Pourquoi la formule fonctionne}

On peut découper un triangle à une extrémité du parallélogramme et le déplacer à l'autre extrémité.

On obtient un rectangle de même base et de même hauteur.

L'aire est donc conservée :
\[
A=b\times h.
\]

Ce raisonnement par découpage met en évidence un invariant d'aire.

\section{25. Figures complexes}

Pour calculer l'aire d'une figure composée :
\begin{itemize}
\item on peut la découper en figures usuelles ;
\item ou compléter une figure plus simple puis soustraire.
\end{itemize}

Les conversions d'unités doivent être effectuées avant les additions ou soustractions d'aires.

\section{26. Exemple développé}

Une figure est composée d'un parallélogramme d'aire :
\[
36\ \text{cm}^2
\]
et d'un triangle d'aire :
\[
14\ \text{cm}^2.
\]

S'ils ne se recouvrent pas, l'aire totale vaut :
\[
36+14=50\ \text{cm}^2.
\]

Donc :
\[
\boxed{50\ \text{cm}^2}.
\]

\section{27. Méthode de construction}

\begin{methode}
Pour une construction géométrique :
\begin{enumerate}
\item lister les données ;
\item choisir l'instrument adapté à chaque donnée ;
\item construire progressivement ;
\item conserver les traits utiles ;
\item vérifier les longueurs et angles obtenus ;
\item coder la figure.
\end{enumerate}
\end{methode}

\section{28. Méthode de démonstration}

\begin{methode}
Pour prouver qu'un quadrilatère est un parallélogramme :
\begin{enumerate}
\item identifier les données disponibles ;
\item choisir une propriété caractéristique adaptée ;
\item établir ses conditions ;
\item citer la propriété ;
\item conclure.
\end{enumerate}
\end{methode}

\section{29. Erreurs fréquentes}

\begin{itemize}
\item Confondre hauteur, médiane et médiatrice.
\item Penser qu'une hauteur est toujours à l'intérieur du triangle.
\item Utiliser la moitié d'une base sans justification.
\item Confondre aire du triangle et aire du parallélogramme.
\item Croire qu'un quadrilatère « ressemble » à un parallélogramme suffit pour conclure.
\item Utiliser une propriété caractéristique sans vérifier toutes ses conditions.
\item Confondre rectangle, losange et carré.
\item Oublier les unités carrées.
\end{itemize}

\section{À retenir}

\begin{aretenir}
Dans tout triangle :
\[
\boxed{\widehat A+\widehat B+\widehat C=180^\circ}.
\]

Son aire vaut :
\[
\boxed{A=\dfrac{b\times h}{2}}.
\]

Les hauteurs sont concourantes.

Une médiane partage le triangle en deux triangles de même aire.

Un parallélogramme possède des côtés opposés parallèles et ses diagonales se coupent en leur milieu.

Son aire vaut :
\[
\boxed{A=b\times h}.
\]

Rectangle, losange et carré sont des parallélogrammes particuliers possédant des propriétés supplémentaires.
\end{aretenir}
